\documentclass[../../main.tex]{subfiles}% 注意这里的文件路径不能用 ./main.tex ,否则用latexmk编译子文件会报错
\graphicspath{{\subfix{./image/}}} % 指定图片目录,后续可以直接使用图片文件名
% 注意这里的文件路径不能用 ../../image/ ,否则用latexmk编译子文件会报错

% 例如:
% \begin{figure}[H]
% \centering
% \includegraphics[scale=0.3]{图.png}
% \caption{}
% \label{figure:图}
% \end{figure}
% 注意:上述\label{}一定要放在\caption{}之后,否则引用图片序号会只会显示??.

\begin{document}

\section{可分扩张$\,\,$本原元素}

\begin{definition}
设$K$是域$F$的扩域，$\alpha\in K$是$F$上的代数元。若$\mathrm{Irr}(\alpha,F)$可分，则称$\alpha$是$F$上的\textbf{可分元素}；若$\mathrm{Irr}(\alpha,F)$不可分，则称$\alpha$是$F$上的\textbf{不可分元素}。
$\alpha$为可分元素当且仅当$\mathrm{Irr}(\alpha,F)'\neq0$。
\end{definition}

\begin{definition}
设$K$是域$F$的代数扩张。若$K$中每个元素都是$F$上的可分元素，则称$K$是$F$的\textbf{可分扩张}，否则称为\textbf{不可分扩张}。
\end{definition}

\begin{example}
完备域$F$的任何代数扩张$K$都是$F$的可分扩张。
\end{example}
\begin{proof}
任取$\alpha\in K$，$\mathrm{Irr}(\alpha,F)\in F[x]$为不可约多项式。由$F$为完备域，$\alpha$是$F$上的可分元素，故$K$是$F$的可分扩张。

\end{proof}

\begin{example}
有限域$F$的任何代数扩张$K$都是可分扩张。
\end{example}
\begin{proof}
有限域是完备域，故结论成立。

\end{proof}

\begin{example}
若域$F$的特征为$0$，则$F$的任何代数扩张都是可分扩张。
\end{example}
\begin{proof}
特征为$0$的域是完备域，故结论成立。

\end{proof}

\begin{example}
设$p$是素数，$\mathbb{Z}_p(t)$是$\mathbb{Z}_p$的单超越扩张，$\alpha$是$x^p-t\in\mathbb{Z}_p(t)[x]$的根，则$\mathbb{Z}_p(t)$的代数扩张$\mathbb{Z}_p(t)(\alpha)$是$\mathbb{Z}_p(t)$的不可分扩张。
\end{example}
\begin{proof}
$x^p-t$是不可分多项式，故$\alpha$为不可分元素，因此该扩张为不可分扩张。

\end{proof}

\begin{theorem}
域$F$上可分多项式$f(x)$的分裂域$K$是$F$的可分扩张。
\end{theorem}
\begin{proof}
$f(x)$为可分多项式，故其每个不可约因子无重根。$K$是$f(x)$的分裂域，任取$\alpha\in K$，$\mathrm{Irr}(\alpha,F)$的根全在$K$中，故$K$也是$\mathrm{Irr}(\alpha,F)f(x)$的分裂域。$\mathrm{Irr}(\alpha,F)f(x)$的每个不可约多项式无重根，故$\mathrm{Irr}(\alpha,F)$是$F$上的可分多项式，即$\alpha$是$F$上的可分元素。因此$K$是$F$的可分扩张。

\end{proof}

\begin{corollary}
设$F(\alpha_1,\alpha_2,\dots,\alpha_n)$是域$F$的代数扩张，且$\alpha_1,\alpha_2,\dots,\alpha_n$是$F$上的可分元素，则该扩张为$F$上的可分扩张。
\end{corollary}
\begin{proof}
\begin{align*}
f(x)=\prod\limits_{i=1}^{n}\mathrm{Irr}(\alpha_i,F),
\end{align*}
则$f(x)$是$F$上可分多项式，其分裂域$K$是$F$上的可分扩张。对任意$\alpha\in F(\alpha_1,\alpha_2,\dots,\alpha_n)\subseteq K$，$\alpha$是$F$上的可分元素，故$F(\alpha_1,\alpha_2,\dots,\alpha_n)$是$F$上的可分扩张。

\end{proof}

\begin{definition}
设$K$是域$F$的有限扩张。若存在$\alpha\in K$，使得$K=F(\alpha)$，则称$\alpha$是$K$对于$F$的\textbf{本原元素}。
\end{definition}

\begin{theorem}
设$K=F(\alpha_1,\alpha_2,\dots,\alpha_n)$是域$F$的代数扩张，且$\alpha_1,\alpha_2,\dots,\alpha_n$是$F$上的可分元素，则$K$中有本原元素，即$K$是$F$的单代数扩张。
\end{theorem}
\begin{proof}
分两种情况讨论：
\begin{enumerate}[(1)]
\item 若$F$是有限域，则$K$也是有限域。$K^*=K\setminus\{0\}$是循环群，存在$\alpha\in K^*$使得$K^*=\langle\alpha\rangle$，故$K=F(\alpha)$，$\alpha$是本原元素。
\item 若$F$是无限域，对添加元素个数$n$用数学归纳法。
当$n=1$时，$\alpha_1$即为本原元素。
设$n-1$时定理成立，即存在$\beta\in F(\alpha_1,\alpha_2,\dots,\alpha_{n-1})$，使得$F(\alpha_1,\alpha_2,\dots,\alpha_{n-1})=F(\beta)$，于是$K=F(\beta,\alpha_n)$。

设$E$为$\mathrm{Irr}(\beta,F)\cdot\mathrm{Irr}(\alpha_n,F)$的分裂域，在$E[x]$中有分解
\begin{align*}
\mathrm{Irr}(\beta,F)&=(x-\beta)(x-\beta_2)\cdots(x-\beta_s),\\
\mathrm{Irr}(\alpha_n,F)&=(x-\alpha_n)(x-\alpha_{n2})\cdots(x-\alpha_{nt}).
\end{align*}
因$\alpha_n$可分，故$i\neq j$时，$\alpha_{ni}\neq\alpha_{nj}$。记$\beta=\beta_1$，集合$\left\{\frac{\beta_k-\beta}{\alpha_n-\alpha_{nj}}\right\}$为有限集。由$F$无限，可取$c\in F$使得
\begin{align*}
c\neq\frac{\beta_k-\beta}{\alpha_n-\alpha_{nj}}.
\end{align*}

令$\theta=\beta+c\alpha_n$，构造多项式
\begin{align*}
f(x)=((\theta-cx)-\beta)((\theta-cx)-\beta_2)\cdots((\theta-cx)-\beta_s).
\end{align*}
由$\theta-c\alpha_n=\beta$得$f(\alpha_n)=0$；对任意$2\leqslant j\leqslant t$，有$\theta-c\alpha_{nj}-\beta_k=\beta+c\alpha_n-c\alpha_{nj}-\beta_k\neq0$，故$f(\alpha_{nj})\neq0$。
因此$(f(x),\mathrm{Irr}(\alpha_n,F))=x-\alpha_n$。
结合$f(x),\mathrm{Irr}(\alpha_n,F)\in F(\theta)[x]$，得$\alpha_n\in F(\theta)$，进而$\beta=\theta-c\alpha_n\in F(\theta)$。于是$F(\theta)\subseteq F(\beta,\alpha_n)\subseteq F(\theta)$，即$F(\theta)=F(\alpha_1,\alpha_2,\dots,\alpha_n)$，$\theta$是本原元素。
\end{enumerate}

\end{proof}

\begin{corollary}
一个域的有限可分扩张一定是单代数扩张。
\end{corollary}

\begin{lemma}
设域$F$的特征为$p\neq0$，则$F$的单代数扩张$F(\alpha)$为可分扩张的充分必要条件为$F(\alpha)=F(\alpha^p)$。
\end{lemma}
\begin{proof}
由前述推论，$F(\alpha)$为可分扩张当且仅当$\alpha$为可分元素，故只需证$\alpha$可分当且仅当$F(\alpha)=F(\alpha^p)$。

若$\alpha$是$F$上的可分元素，则$\alpha$是$F(\alpha^p)$上的可分元素。记$f(x)=x^p-\alpha^p\in F(\alpha^p)[x]$，$\alpha$是$f(x)$的根。设$E$是$f(x)$的分裂域，则$f(x)=(x-\alpha)^p$，于是$\mathrm{Irr}(\alpha,F(\alpha^p))=(x-\alpha)^m$，$1\leqslant m\leqslant p$。
由$\alpha$在$F(\alpha^p)$上可分，得$m=1$，即$\alpha\in F(\alpha^p)$，故$F(\alpha)=F(\alpha^p)$。

反之，若$\alpha$在$F$上不可分，则存在$F[x]$中可分的不可约多项式$g(x)$，使得$\mathrm{Irr}(\alpha,F)=g(x^p)$。于是$\alpha^p$是$g(x)$的根，$\mathrm{Irr}(\alpha^p,F)=g(x)$。由$\deg(\alpha,F)>\deg(\alpha^p,F)$，得$F(\alpha)\supsetneq F(\alpha^p)$。

\end{proof}

\begin{lemma}
设$F(\alpha,\beta)$是域$F$的代数扩张，$\alpha$是$F$上的可分元素，$\beta$是$F(\alpha)$上的可分元素，则$F(\alpha,\beta)$是$F$的可分扩张。
\end{lemma}
\begin{proof}
若$\mathrm{ch}F=0$，结论显然成立。下设$\mathrm{ch}F=p\neq0$，只需证$\beta$是$F$上的可分元素，即证$F(\beta)=F(\beta^p)$。

由域塔扩张次数公式：
\begin{align*}
[F(\alpha,\beta):F(\beta^p)]=[F(\alpha,\beta):F(\beta)]\cdot[F(\beta):F(\beta^p)],
\end{align*}
故只需证$[F(\alpha,\beta):F(\beta^p)]=[F(\alpha,\beta):F(\beta)]$。

$\alpha$在$F$上可分，故在$F(\beta),F(\beta^p)$上均可分。记$g(x)=\mathrm{Irr}(\alpha,F(\beta))$，则$g(x)^p\in F(\beta^p)[x]\subseteq F(\beta)[x]$。记$h(x)=\mathrm{Irr}(\alpha,F(\beta^p))$，则$g(x)\mid h(x)$，$h(x)\mid g(x)^p$。
$h(x)$与$g(x)$无重根，故$g(x)=h(x)$，因此
\begin{align*}
[F(\alpha,\beta^p):F(\beta^p)]=[F(\alpha,\beta):F(\beta)].
\end{align*}

$\beta$在$F(\alpha)$上可分，由前述引理得$F(\alpha,\beta^p)=F(\alpha)(\beta^p)=F(\alpha)(\beta)=F(\alpha,\beta)$，故$\beta$在$F$上可分，$F(\alpha,\beta)$是$F$的可分扩张。

\end{proof}

\begin{theorem}
设$E$是域$F$的代数扩张，$K$是中间域，即$F\subseteq K\subseteq E$，则$E$是$F$的可分扩张当且仅当$K$是$F$的可分扩张，且$E$是$K$的可分扩张。
\end{theorem}
\begin{proof}
必要性显然成立，下证充分性。
任取$\alpha\in E$，由$\alpha$是$K$上的可分元素，记
\begin{align*}
\mathrm{Irr}(\alpha,K)=x^n+a_{n-1}x^{n-1}+\dots+a_0,
\end{align*}
其为$K$上可分不可约多项式，故$\alpha$是$F(a_0,a_1,\dots,a_{n-1})$上的可分元素。

$K$是$F$的可分扩张，故$a_0,a_1,\dots,a_{n-1}$是$F$上的可分元素。由本原元定理，存在$\theta$使得$F(a_0,a_1,\dots,a_{n-1})=F(\theta)$。
于是$\alpha$是$F(\theta)$上的可分元素，$\theta$是$F$上可分元素，故$F(\theta,\alpha)$是$F$的可分扩张，即$\alpha$是$F$上的可分元素。因此$E$是$F$的可分扩张。

\end{proof}



\end{document}